\section{Introduction}
\label{sec:intro}

Legislators routinely reach for scientific authority. A bill is ``backed by research,'' a problem
is established with a statistic, an intervention is said to be ``proven'' to work. In technically
complex domains---and few are more complex or fast-moving than artificial intelligence (AI)---these
appeals to evidence are both more common and harder to evaluate. Yet a basic question has remained
effectively unmeasurable at scale: when a legislator invokes science to support a claim, does the
scientific record actually support that claim? Answering it would let us move the study of research
use in policy \citep{weiss1979many,oliver2022works} from case studies and surveys toward systematic,
population-level measurement.

The obstacle is not data but \emph{validation}. Legislators' public communications are abundant and
machine-readable, and modern language models can extract the empirical claims they contain and
retrieve plausibly relevant scientific literature. The hard step is the last one: deciding whether a
retrieved study \emph{supports}, \emph{refutes}, or is \emph{silent} on a given claim. This judgment
has traditionally required trained human coders, whose cost and scarcity cap the size and scope of
any such study and whose reliability must itself be established. As a practical matter, the
``does-the-science-support-the-claim'' step is where projects of this kind stall.

We show that this step can be performed \emph{autonomously}, with no human research assistants and no
proprietary application programming interfaces (APIs), while preserving statistical validity. Our
design rests on three ideas from the recent literature. First, high-capacity large language models
(LLMs) now match or exceed trained crowd coders on relevance and stance tasks
\citep{gilardi2023chatgpt,ziems2024can}, so an LLM can supply gold-standard labels. Second, because
any machine labels are imperfect and naively substituting them for human labels biases downstream
estimates \citep{egami2023dsl}, we combine a small set of high-quality LLM-judge labels with a large
set of cheap, independent classifier labels using prediction-powered inference (PPI)
\citep{angelopoulos2023ppi} and design-based supervised learning (DSL) \citep{egami2023dsl}. Third,
because any single judge can be biased \citep{zheng2023judging}, we pair a generative LLM judge with
discriminative natural-language-inference (NLI) classifiers that fail in different ways, and we report
their agreement explicitly, in the spirit of the alternative-annotator test \citep{calderon2025alttest}.

We apply the system to a corpus of 7{,}139 California State Assembly press releases, 303 of which
discuss AI, and to a parallel pilot in K--12 education policy, and we report two results cautiously.
\emph{First, methodologically,} our \emph{zero-shot} NLI configuration agrees with the LLM judge only
weakly ($\kappa \approx 0.15$--$0.21$) and labels most pairs ``silent,'' because strict entailment
between an abstract and a press-release claim is usually ``neutral'' even when the study substantiates
the claim's substance. Anchoring on the LLM judge and applying PPI moves the estimated support rate from
a naive $0.21$ (Democrats) / $0.11$ (Republicans) to $0.71$ / $0.54$. We resist the tempting reading that
``the LLM judge recovers the truth.'' A correction this large ($\sim$3$\times$) is atypical when a
surrogate is competent, and is equally consistent with a misconfigured NLI baseline (zero-shot entailment
is a poor proxy for scientific support; serious claim verification uses fine-tuning and rationale
selection \citep{wadden2020fact}) and with the LLM judge \emph{over}-calling support. Because the claims
and references were produced by LLMs and our judge is an LLM, \emph{self-preference}---the documented
tendency of LLM evaluators to favor model-generated text \citep{panickssery2024llm}---could inflate the
support rate, and we have no human anchor to rule this out. \emph{Second, substantively and
provisionally,} under our judge both parties' AI claims are supported at broadly similar (and fairly
high) rates, so the large partisan gap implied by the naive NLI estimate does not survive; what is robust
is that Democrats discuss AI far more often and more intensively than Republicans, not any difference in
the support for their claims.

Our contributions are deliberately modest. (1) We provide a reusable, end-to-end, fully autonomous
\emph{template}---claim extraction, evidence retrieval, and support assessment---that runs with no human
coders and no paid APIs, released as open-source software, that others can apply and, crucially, validate
against a human standard. (2) We show that a popular shortcut---scoring scientific support with off-the-shelf
zero-shot NLI---can be badly miscalibrated for press-release claims, a cautionary result for the growing
practice of measuring political text with such tools. (3) We are explicit about what an autonomous,
single-model, self-judging pipeline \emph{cannot} establish on its own (validity against truth), and we
specify what a valid version would require. We do \emph{not} claim a validated estimate of evidence use;
the substantive patterns we report are provisional and conditional on an unvalidated judge. The remainder of the paper describes the related literatures (\S\ref{sec:related}),
the data (\S\ref{sec:data}), our methods (\S\ref{sec:methods}), and our results
(\S\ref{sec:results}), before discussing implications and limitations
(\S\ref{sec:discussion}--\S\ref{sec:conclusion}).
